\par\noindent\textbf{(1)} 短天线两端累积电荷满足连续性 \(i\omega Q_0=I_0\)，偶极矩幅 \(p_0=Q_0l=I_0l/\omega\)（相位依时谐约定）。把天线当作赫兹短偶极子必须满足沿长度电流近似均匀，否则有效长度应由电流分布积分。
\par\noindent\textbf{(2)} 电偶极总平均辐射功率为 \(P_{\rm rad}=\omega^4|p_0|^2/(12\pi\varepsilon_0c^3)\)。代入 \(p_0=I_0l/\omega\) 得 \(P_{\rm rad}=I_0^2\omega^2l^2/(12\pi\varepsilon_0c^3)=I_0^2\eta_0k^2l^2/(12\pi)\)，\(k=\omega/c\)。
\par\noindent\textbf{(3)} 以电流峰值约定 \(P_{\rm rad}=I_0^2R_{\rm rad}/2\)，得 \(R_{\rm rad}=\eta_0(k l)^2/(6\pi)=80\pi^2(l/\lambda)^2\,\mathrm{ohm}\)（真空 \(\eta_0\approx120\pi\,\mathrm{ohm}\)）。因此天线尺寸减半时辐射电阻降为四分之一。
\par\noindent\textbf{(4)} 若馈线等效电阻为 \(R_{\rm ohm}\)，在谐振阻抗匹配近似下辐射效率 \(\eta=R_{\rm rad}/(R_{\rm rad}+R_{\rm ohm})\)。\(\eta\) 不能通过忽略 \(R_{\rm ohm}\) 就推断很小天线仍高效；大部分输入功率可能转为热。
